% Options for packages loaded elsewhere
\PassOptionsToPackage{unicode}{hyperref}
\PassOptionsToPackage{hyphens}{url}
\documentclass[
  11pt,
]{article}
\usepackage{xcolor}
\usepackage[left=3cm,right=3cm,top=2.5cm,bottom=2.5cm]{geometry}
\usepackage{amsmath,amssymb}
\setcounter{secnumdepth}{5}
\usepackage{iftex}
\ifPDFTeX
  \usepackage[T1]{fontenc}
  \usepackage[utf8]{inputenc}
  \usepackage{textcomp} % provide euro and other symbols
\else % if luatex or xetex
  \usepackage{unicode-math} % this also loads fontspec
  \defaultfontfeatures{Scale=MatchLowercase}
  \defaultfontfeatures[\rmfamily]{Ligatures=TeX,Scale=1}
\fi
\usepackage{lmodern}
\ifPDFTeX\else
  % xetex/luatex font selection
\fi
% Use upquote if available, for straight quotes in verbatim environments
\IfFileExists{upquote.sty}{\usepackage{upquote}}{}
\IfFileExists{microtype.sty}{% use microtype if available
  \usepackage[]{microtype}
  \UseMicrotypeSet[protrusion]{basicmath} % disable protrusion for tt fonts
}{}
\usepackage{setspace}
\makeatletter
\@ifundefined{KOMAClassName}{% if non-KOMA class
  \IfFileExists{parskip.sty}{%
    \usepackage{parskip}
  }{% else
    \setlength{\parindent}{0pt}
    \setlength{\parskip}{6pt plus 2pt minus 1pt}}
}{% if KOMA class
  \KOMAoptions{parskip=half}}
\makeatother
\usepackage{longtable,booktabs,array}
\usepackage{caption}
\captionsetup[table]{skip=6pt}
\newcounter{none} % for unnumbered tables
\usepackage{calc} % for calculating minipage widths
% Correct order of tables after \paragraph or \subparagraph
\usepackage{etoolbox}
\makeatletter
\patchcmd\longtable{\par}{\if@noskipsec\mbox{}\fi\par}{}{}
\makeatother
% Allow footnotes in longtable head/foot
\IfFileExists{footnotehyper.sty}{\usepackage{footnotehyper}}{\usepackage{footnote}}
\makesavenoteenv{longtable}
\usepackage{graphicx}
\makeatletter
\newsavebox\pandoc@box
\newcommand*\pandocbounded[1]{% scales image to fit in text height/width
  \sbox\pandoc@box{#1}%
  \Gscale@div\@tempa{\textheight}{\dimexpr\ht\pandoc@box+\dp\pandoc@box\relax}%
  \Gscale@div\@tempb{\linewidth}{\wd\pandoc@box}%
  \ifdim\@tempb\p@<\@tempa\p@\let\@tempa\@tempb\fi% select the smaller of both
  \ifdim\@tempa\p@<\p@\scalebox{\@tempa}{\usebox\pandoc@box}%
  \else\usebox{\pandoc@box}%
  \fi%
}
% Set default figure placement to htbp
\def\fps@figure{htbp}
\makeatother
\setlength{\emergencystretch}{3em} % prevent overfull lines
\providecommand{\tightlist}{%
  \setlength{\itemsep}{0pt}\setlength{\parskip}{0pt}}
%% tools/stamp.tex
%% Provenance footer for PDFs rendered from this project.
%% Vendored by zzcollab; this file is not project-specific. The
%% footer prints the source document and the time of rendering.
%% The git version is defined here too, and so is available to a
%% project that wants it on the page, but it is left off the
%% default footer: it already names the staged copy in share/ and
%% appears in its manifest row. All values are supplied at render
%% time by tools/stamp-render.R, which writes a small generated
%% file that redefines the macros below. The \providecommand
%% defaults keep a bare LaTeX compile working when the document is
%% built without the wrapper.
\usepackage{fancyhdr}
\usepackage{xcolor}
\providecommand{\stampsource}{\detokenize{(source not recorded)}}
\providecommand{\stamptime}{(time not recorded)}
\providecommand{\stampversion}{\detokenize{(version not recorded)}}
\setlength{\footskip}{40pt}
\newcommand{\stampfootcontent}{%
  \footnotesize\color{gray}%
  \parbox[b]{\textwidth}{%
    \stamptime\hfill\thepage\\[2pt]%
    \ttfamily\raggedright\stampsource}}
\pagestyle{fancy}
\fancyhf{}
\fancyfoot[C]{\stampfootcontent}
\renewcommand{\headrulewidth}{0pt}
\renewcommand{\footrulewidth}{0.4pt}
%% The title page uses the `plain` page style; redefine it so the
%% stamp appears there too.
\fancypagestyle{plain}{%
  \fancyhf{}%
  \fancyfoot[C]{\stampfootcontent}%
  \renewcommand{\headrulewidth}{0pt}%
  \renewcommand{\footrulewidth}{0.4pt}}
\renewcommand{\stampsource}{\textasciitilde{}/\allowbreak{}prj/\allowbreak{}res/\allowbreak{}04-runin-power-analysis/\allowbreak{}runinpower/\allowbreak{}analysis/\allowbreak{}report/\allowbreak{}01-runin-power/\allowbreak{}bullets.Rmd}
\renewcommand{\stamptime}{2026-08-16 16:12 PDT}
\renewcommand{\stampversion}{\detokenize{276ac7f-wip}}
\usepackage{amsmath}
\usepackage{amssymb}
\usepackage{booktabs}
\usepackage{longtable}
\usepackage{float}
\usepackage[mathlines]{lineno}
\linenumbers
\DeclareMathOperator{\Var}{Var}
\DeclareMathOperator{\Cov}{Cov}
\DeclareMathOperator{\tr}{tr}
\usepackage{bookmark}
\IfFileExists{xurl.sty}{\usepackage{xurl}}{} % add URL line breaks if available
\urlstyle{same}
% fallback for those not using the hyperref driver hyperxmp:
\makeatletter
\@ifundefined{xmpquote}{\newcommand{\xmpquote}[1]{#1}}{}
\makeatother
\hypersetup{
  pdftitle={Structured summary: \textbar{} Power Analysis for Longitudinal Clinical Trials with Run-In and Common Close Observations},
  pdfauthor={Ronald G. Thomas, Ph.D.},
  hidelinks,
  pdfcreator={LaTeX via pandoc}}

\title{Structured summary: \textbar{} Power Analysis for Longitudinal Clinical Trials with Run-In and Common Close Observations}
\author{Ronald G. Thomas, Ph.D.\thanks{Wertheim School of Public Health, UCSD; ORCID: 0000-0003-1686-4965}}
\date{August 16, 2026}

\begin{document}
\maketitle

{
\setcounter{tocdepth}{2}
\tableofcontents
}
\setstretch{1.1}
\emph{This document is a structured, section-by-section summary of the key points argued in the companion manuscript, ``| Power Analysis for Longitudinal Clinical Trials with Run-In and Common Close Observations.'' It is provided as a supplementary reading aid and does not stand on its own; all notation, derivations, simulation detail, and citations are in the main manuscript.}

\bigskip\noindent

\rule{\textwidth}{0.4pt}\bigskip

\section{Introduction}\label{introduction}

\textbf{Beyond its traditional roles, the run-in period has a distinct methodological value in rate-of-change trials: estimating individual baseline trajectories to raise power.}

\begin{itemize}\setlength{\itemsep}{2pt}\setlength{\parskip}{0pt}
\item Conventional run-in uses include identifying compliant participants and washing out prior treatments.
\item When the outcome is a rate of change, run-in observations estimate each participant's baseline trajectory.
\item Incorporating those trajectories into the analysis model substantially increases statistical power.
\end{itemize}

\textbf{Frost and colleagues established the mixed-model foundation: run-in observations help most when slope variability is large relative to measurement error.}

\begin{itemize}\setlength{\itemsep}{2pt}\setlength{\parskip}{0pt}
\item The key insight is that pre-randomization observations inform each participant's underlying trajectory.
\item This reduces the residual variance of the treatment-effect estimator, and the gain grows with the slope-to-noise ratio.
\item In Alzheimer's neuroimaging trials, where that ratio is favorable, they reported sample-size reductions of 30 to 45 percent.
\end{itemize}

\textbf{Wang and colleagues extended the framework with two-period models that jointly model run-in and post-randomization data rather than treating run-in as a covariate.}

\begin{itemize}\setlength{\itemsep}{2pt}\setlength{\parskip}{0pt}
\item Joint modelling of both periods yields additional power gains of up to 15 percent.
\item It permits unequal randomization ratios without power loss.
\item Run-in observations contribute placebo-equivalent information for all participants.
\end{itemize}

\textbf{A parallel literature gives closed-form variance and sample-size formulae for the random-coefficient model without a run-in phase, which form the baseline this paper extends.}

\begin{itemize}\setlength{\itemsep}{2pt}\setlength{\parskip}{0pt}
\item Lu and colleagues derived the cLDA estimator variance, its link to ANCOVA, and the conditions under which the two coincide.
\item Hu and colleagues generalized these to random-coefficient models with monotone missing data and variable follow-up, illustrated with ADNI data.
\item These results cover only the $J_0 = 0$ case and serve as the baseline extended here.
\end{itemize}

\textbf{The CRAN longpower package is the reference implementation of the baseline formulae, and our general result reduces to its routines at the appropriate boundaries.}

\begin{itemize}\setlength{\itemsep}{2pt}\setlength{\parskip}{0pt}
\item longpower provides edland, diggle, and liu.liang routines for slope-comparison trials.
\item The Stata package slopepower offers a complementary simulation-based route to the same class of calculations.
\item Our formula reduces to the Ard-Edland slope variance at $J_0 = J_2 = 0$, and to the Diggle compound-symmetry form when also $\sigma_b^2 = 0$.
\item The generalization thus embeds the existing longpower formulae as the baseline against which run-in and common-close gains are measured.
\end{itemize}

\textbf{No prior closed-form power formula handles an arbitrary number of run-in observations together with a separate common-close phase, the gap this paper fills.}

\begin{itemize}\setlength{\itemsep}{2pt}\setlength{\parskip}{0pt}
\item The common-close phase observes all participants off-treatment after the randomized period.
\item It supplies extra information about individual trajectories at no cost to the treatment comparison.
\item Such post-treatment periods appear in delayed-start and randomized-withdrawal designs, but their GLS efficiency contribution has not been formally characterized.
\end{itemize}

\textbf{The paper presents three nested results, all derived from the GLS estimator via the Woodbury identity.}

\begin{itemize}\setlength{\itemsep}{2pt}\setlength{\parskip}{0pt}
\item A closed-form variance for exactly two run-in observations (Section 3.2), then its generalization to arbitrary $J_0$ (Section 3.3).
\item A further extension adding $J_2$ common-close observations (Section 3.4).
\item All expressions are in terms of $\sigma^2$, $\sigma_b^2$, and spacing $t$, with numerical examples from Alzheimer's neuroimaging studies (Section 5).
\end{itemize}

\subsection{General formulation}\label{general-formulation}

\textbf{The model accommodates a two-arm longitudinal trial whose observation schedule spans three distinct phases.}

\begin{itemize}\setlength{\itemsep}{2pt}\setlength{\parskip}{0pt}
\item $N$ participants are observed at equally spaced time points.
\item The phases are run-in (pre-randomization), treatment (post-randomization), and an optional common close (post-treatment).
\item $Y_{ij}$ denotes the outcome for participant $i$ at time point $j$.
\end{itemize}

\textbf{The outcome model combines a random intercept and random slope with a three-component within-subject variance structure following Frost and colleagues.}

\begin{itemize}\setlength{\itemsep}{2pt}\setlength{\parskip}{0pt}
\item $a_i$ and $b_i$ are the random intercept and slope, while $u_{ij}$ and $\varepsilon_{ij}$ are the between-visit and between-scan residuals.
\item Both residual terms are independent across visits, so under a single scan per visit they enter every formula only through their sum.
\item We therefore define the total within-visit residual $w_{ij} = u_{ij} + \varepsilon_{ij}$ with variance $\sigma^2$.
\end{itemize}

\textbf{The two within-visit components are identifiable only with replicate scans, so single-scan designs depend solely on their sum.}

\begin{itemize}\setlength{\itemsep}{2pt}\setlength{\parskip}{0pt}
\item Separating $\sigma_u^2$ from $\sigma_\varepsilon^2$ requires a visit carrying replicate scans.
\item For the single-scan designs treated here, only the sum $\sigma^2$ enters the treatment-effect variance.
\item This reduction simplifies all subsequent variance expressions.
\end{itemize}

\textbf{Three fixed-effect slope parameters describe the run-in, placebo post-randomization, and treatment trajectories under equally spaced observations.}

\begin{itemize}\setlength{\itemsep}{2pt}\setlength{\parskip}{0pt}
\item $\delta$ is the run-in rate of change common to all participants.
\item $\beta$ is the placebo post-randomization rate of change, and $\gamma$ is the treatment effect on that rate.
\item Observations are equally spaced with interval $t$, so $t_j = jt$.
\end{itemize}

\textbf{During the common close both arms revert to a shared off-treatment slope, modelled by deactivating the treatment indicator in that phase.}

\begin{itemize}\setlength{\itemsep}{2pt}\setlength{\parskip}{0pt}
\item For $j > J_1$ all participants are off-treatment with slope $\beta + b_i$.
\item This is implemented by setting $g_i = 0$ for every participant during the common close.
\item Equivalently, the treatment indicator is replaced with $g_i \cdot I(1 \le j \le J_1)$ in the model.
\end{itemize}

\subsection{Change-score parameterization}\label{change-score-parameterization}

\textbf{Differencing relative to baseline removes the intercept and aligns the analysis with constrained longitudinal data analysis under compound symmetry.}

\begin{itemize}\setlength{\itemsep}{2pt}\setlength{\parskip}{0pt}
\item The change score eliminates $\alpha$ and the random intercept $a_i$, leaving only two variance components.
\item It is algebraically equivalent to cLDA, so the derived formulas apply to both analyses of run-in designs.
\item By dropping $\sigma_a^2$ it avoids the anticonservative sample sizes that arise from ignoring intercept-slope correlation.
\end{itemize}

\subsection{Covariance structure}\label{covariance-structure}

\textbf{The shared baseline residual induces compound-symmetry correlation among change scores, so the residual covariance depends only on the summed within-visit variance.}

\begin{itemize}\setlength{\itemsep}{2pt}\setlength{\parskip}{0pt}
\item The within-visit residuals are independent across visits with common variance $\sigma^2$.
\item Differencing against the shared baseline $w_{i0}$ produces positive correlation, giving $R = \sigma^2(\mathbf{I}_J + \mathbf{1}_J \mathbf{1}_J')$.
\item Both within-visit components enter $R$ identically, so only their sum is identified from single-scan designs.
\end{itemize}

\subsection{GLS estimation and the Woodbury identity}\label{gls-estimation-and-the-woodbury-identity}

\textbf{The estimand of interest is the treatment-effect variance, a single diagonal element of the inverted GLS information matrix.}

\begin{itemize}\setlength{\itemsep}{2pt}\setlength{\parskip}{0pt}
\item The GLS estimator has covariance $(X' \Sigma^{-1} X)^{-1}$.
\item $X$ is the fixed-effects design matrix and $\mathbf{c}$ is the stacked change-score vector.
\item The target $\mathop{\mathrm{Var}}(\hat{\gamma})$ is the diagonal entry corresponding to $\gamma$.
\end{itemize}

\subsubsection{Per-participant information, compact form}\label{per-participant-information-compact-form}

\textbf{The additive per-participant information form has two properties exploited later: a scalar rank-one correction and accommodation of heterogeneous designs.}

\begin{itemize}\setlength{\itemsep}{2pt}\setlength{\parskip}{0pt}
\item The total information is the sum of per-participant contributions, with covariance given by its inverse.
\item Under a single random slope, $\mathcal{J}_i$ is a scalar, making the correction a rank-one adjustment of the OLS-style information.
\item Because contributions add, varying $J_0$, $J_2$, or dropout patterns are handled by summing over strata.
\end{itemize}

\subsection{\texorpdfstring{Case 1: Standard design, no run-in (\(J_0=0\), \(J_1=2\), \(J_2=0\))}{Case 1: Standard design, no run-in (J\_0=0, J\_1=2, J\_2=0)}}\label{case-1-standard-design-no-run-in-j_00-j_12-j_20}

\textbf{The no-run-in case reproduces established results and serves as a verification baseline for the framework.}

\begin{itemize}\setlength{\itemsep}{2pt}\setlength{\parskip}{0pt}
\item It reproduces Frost and colleagues, Section 3.1.1.
\item It is the $J_0 = 0$ baseline of the random-coefficient power formulae of Lu and Hu and colleagues.
\item It thereby provides a verification checkpoint for the present derivation.
\end{itemize}

\textbf{Without a run-in phase the design reduces to two post-randomization change scores and a two-parameter fixed-effects model.}

\begin{itemize}\setlength{\itemsep}{2pt}\setlength{\parskip}{0pt}
\item There are $J = 2$ change scores at post-randomization times $t$ and $2t$.
\item The absence of run-in removes the $\delta$ parameter from the model.
\item The remaining fixed effects are the placebo slope $\beta$ and the treatment effect $\gamma$.
\end{itemize}

\textbf{The recovered variance matches Frost and supplies the per-group baseline for all run-in comparisons.}

\begin{itemize}\setlength{\itemsep}{2pt}\setlength{\parskip}{0pt}
\item The expression agrees with Frost and colleagues, Section 3.1.1.
\item For $m$ participants per group the variance scales as $(\sigma^2 + 2\sigma_b^2 t^2)/(m t^2)$.
\item This quantity is the reference against which run-in designs are evaluated.
\end{itemize}

\subsection{\texorpdfstring{Case 2: Two run-in observations (\(J_0=2\), \(J_1=1\), \(J_2=0\))}{Case 2: Two run-in observations (J\_0=2, J\_1=1, J\_2=0)}}\label{case-2-two-run-in-observations-j_02-j_11-j_20}

\textbf{Two run-in observations yield three change scores and a design matrix that distinguishes the run-in slope from the post-randomization slopes.}

\begin{itemize}\setlength{\itemsep}{2pt}\setlength{\parskip}{0pt}
\item There are $J = 3$ change scores spanning the pre- and post-randomization times.
\item The fixed-effects model now includes the run-in slope $\delta$ alongside $\beta$ and $\gamma$.
\item The design matrix separates the run-in contribution from the post-randomization contribution.
\end{itemize}

\textbf{Because the random slope acts uniformly across visits, the random-effects design vector collapses to the scaled time distances from baseline.}

\begin{itemize}\setlength{\itemsep}{2pt}\setlength{\parskip}{0pt}
\item The random slope $b_i$ applies to all time points alike.
\item The random-effects design therefore reduces to $Z = t(-2,\, -1,\, 1)'$.
\item The run-in change scores are differences of the pre-randomization outcomes from baseline.
\end{itemize}

\textbf{The treatment-effect variance is the relevant diagonal element of the inverted information matrix, obtained numerically and checked against the closed form.}

\begin{itemize}\setlength{\itemsep}{2pt}\setlength{\parskip}{0pt}
\item $\mathop{\mathrm{Var}}(\hat{\gamma})$ is the $(3,3)$ element of $(X'\Sigma^{-1}X)^{-1}$.
\item It is evaluated numerically in the accompanying R code.
\item The numerical value is verified against the derived closed-form result.
\end{itemize}

\textbf{Completing the cofactor and determinant algebra yields the closed-form treatment-effect variance for the two-run-in design.}

\begin{itemize}\setlength{\itemsep}{2pt}\setlength{\parskip}{0pt}
\item Define the auxiliary quantity $D = \sigma^2 + 5\sigma_b^2 t^2$.
\item The $(3,3)$ cofactor simplifies to $6t^4/(aD)$.
\item Dividing the cofactor by the determinant gives the boxed variance below.
\end{itemize}

\textbf{The run-in variance differs structurally from the Frost baseline by involving the residual variance in both numerator and denominator.}

\begin{itemize}\setlength{\itemsep}{2pt}\setlength{\parskip}{0pt}
\item In the Frost formula $\sigma^2$ appears only additively in the numerator.
\item Here $\sigma^2$ enters both the numerator and the denominator.
\item This reflects the extra information extracted from pre-randomization observations.
\end{itemize}

\subsection{\texorpdfstring{General formula: \(J_0\) run-in observations (\(J_2=0\))}{General formula: J\_0 run-in observations (J\_2=0)}}\label{general-formula-j_0-run-in-observations-j_20}

\textbf{The derivation now generalizes to an arbitrary number of run-in observations with no common close.}

\begin{itemize}\setlength{\itemsep}{2pt}\setlength{\parskip}{0pt}
\item The design has arbitrary $J_0$ run-in and $J_1$ post-randomization observations.
\item The common close is absent, so $J_2 = 0$.
\item The total number of change scores is $J = J_0 + J_1$.
\end{itemize}

\textbf{The design matrix is built from run-in and post-randomization time vectors mapping to the three fixed-effect columns.}

\begin{itemize}\setlength{\itemsep}{2pt}\setlength{\parskip}{0pt}
\item The columns correspond to the parameters $\delta$, $\beta$, and $\gamma$.
\item $\mathbf{d}_{J_0}$ collects the negative run-in times scaled by $t$.
\item $\mathbf{d}_{J_1}$ collects the positive post-randomization times scaled by $t$.
\end{itemize}

\subsubsection{\texorpdfstring{Computing \(X'\Sigma^{-1}X\)}{Computing X\textquotesingle\textbackslash Sigma\^{}\{-1\}X}}\label{computing-xsigma-1x}

\textbf{The Woodbury computation is organized around the sum-of-squares of the run-in and post-randomization time indices.}

\begin{itemize}\setlength{\itemsep}{2pt}\setlength{\parskip}{0pt}
\item The derivation proceeds via the Woodbury identity.
\item $S_{J_0}$ is the sum of squared run-in indices in closed form.
\item $S_{J_1}$ is the corresponding sum over post-randomization indices.
\end{itemize}

\textbf{The structure of the information matrix reflects which parameters both groups inform and the opposing signs of the two phase time-sums.}

\begin{itemize}\setlength{\itemsep}{2pt}\setlength{\parskip}{0pt}
\item Entries informed by both groups carry a factor of 2.
\item The $\gamma$-related entries lack this factor because the placebo $\gamma$ column is zero.
\item The positive $(\delta, \beta)$ cross-term arises from the opposite signs of the run-in and post-randomization time sums.
\end{itemize}

\subsubsection{Numerical implementation}\label{numerical-implementation}

\textbf{For general run-in and post-randomization counts the variance has no compact closed form and must be evaluated as a complex ratio.}

\begin{itemize}\setlength{\itemsep}{2pt}\setlength{\parskip}{0pt}
\item No simple closed-form expression exists for arbitrary $J_0$ and $J_1$.
\item The variance depends on the sum-of-squares and sum quantities together with the dimension $J$.
\item The Woodbury correction interacts with the information matrix to produce algebraically complex expressions.
\end{itemize}

\textbf{A matrix-based R implementation evaluates the exact variance for any configuration, with the closed-form cases as checkpoints.}

\begin{itemize}\setlength{\itemsep}{2pt}\setlength{\parskip}{0pt}
\item The computation uses the matrix Woodbury identity in R code.
\item It returns the exact variance for any $(J_0, J_1, J_2)$ configuration.
\item The $J_0=0$ and $J_0=2$ closed forms serve as verification checkpoints.
\end{itemize}

\subsubsection{Remark: recovery of the Edland and Diggle formulae}\label{remark-recovery-of-the-edland-and-diggle-formulae}

\textbf{Setting the run-in count to zero collapses the information matrix to a two-parameter block and yields the slope-comparison variance.}

\begin{itemize}\setlength{\itemsep}{2pt}\setlength{\parskip}{0pt}
\item With $J_0 = 0$ the $\delta$ column vanishes, leaving a $2 \times 2$ block on $(\beta, \gamma)$.
\item Sherman-Morrison cancellations collapse the Woodbury correction against the remaining structure.
\item The result is the per-arm variance of the treatment-effect estimator for $n$ participants.
\end{itemize}

\textbf{The boundary case recovers the Edland, Diggle, and Frost formulae, confirming that the general result embeds the longpower routines as strict special cases.}

\begin{itemize}\setlength{\itemsep}{2pt}\setlength{\parskip}{0pt}
\item The expression is the Ard-Edland slope-comparison formula implemented in the longpower edland routine.
\item Setting the slope variance to zero gives compound symmetry and recovers the Diggle formula.
\item The Frost formula is the two-visit specialization, so the general result quantifies gains accruing only from run-in or common-close observations.
\end{itemize}

\subsection{\texorpdfstring{General formula with common close (\(J_2 \ge 0\))}{General formula with common close (J\_2 \textbackslash ge 0)}}\label{general-formula-with-common-close-j_2-ge-0}

\textbf{The final extension adds common-close observations after the treatment period to the change-score vector.}

\begin{itemize}\setlength{\itemsep}{2pt}\setlength{\parskip}{0pt}
\item $J_2$ common-close observations follow the treatment period.
\item All participants are off-treatment during this phase.
\item The total dimension becomes $J = J_0 + J_1 + J_2$.
\end{itemize}

\textbf{In the common close the treatment effect persists only as an accumulated offset, fixing the design-matrix entry at the treatment-period value.}

\begin{itemize}\setlength{\itemsep}{2pt}\setlength{\parskip}{0pt}
\item $\gamma$ enters only through the effect accumulated during the treatment period.
\item It does not contribute to the slope during the common close.
\item The $\gamma$ column therefore equals $J_1 t$ for all common-close observations in the treatment group.
\end{itemize}

\textbf{The slope and treatment columns of the design matrix encode the post-randomization times and the capped treatment exposure respectively.}

\begin{itemize}\setlength{\itemsep}{2pt}\setlength{\parskip}{0pt}
\item The $\beta$ column holds the post-randomization times of the placebo slope across all later visits.
\item The $\gamma$ column rises through the treatment period then stays at $J_1$ during the common close.
\item The $\gamma$ column is identically zero for placebo participants.
\end{itemize}

\textbf{The common-close extension reuses the same Woodbury computation with only three quantities altered.}

\begin{itemize}\setlength{\itemsep}{2pt}\setlength{\parskip}{0pt}
\item The dimension $J$ increases to include the common-close visits.
\item The $X$ matrix changes only in its $\gamma$ column.
\item The elements of the random-effects vector $Z$ are extended accordingly.
\end{itemize}

\section{Power and Sample Size Formulas}\label{power-and-sample-size-formulas}

\textbf{A single run-in observation never worsens precision, but the extra slope parameter it introduces can absorb nearly all of the information it contributes.}

\begin{itemize}\setlength{\itemsep}{2pt}\setlength{\parskip}{0pt}
\item Adding one run-in observation introduces the parameter $\delta$, enlarging the fixed-effect dimension from 2 to 3.
\item Across a wide numerical grid, $\mathop{\mathrm{Var}}(\hat\gamma)|_{J_0=1}$ never exceeds $\mathop{\mathrm{Var}}(\hat\gamma)|_{J_0=0}$.
\item At some parameter values the gain is negligible, so the benefit of a single run-in observation should be quantified before extending trial duration.
\end{itemize}

\subsection{Verification of closed-form results}\label{verification-of-closed-form-results}

\textbf{The closed-form expressions are validated against the matrix computation in two representative design configurations.}

\begin{itemize}\setlength{\itemsep}{2pt}\setlength{\parskip}{0pt}
\item The verification confirms agreement between the analytic and matrix-based variances.
\item The Frost standard design uses the two-parameter model with two post-baseline visits.
\item The two-run-in case exercises the three-parameter model.
\end{itemize}

\subsection{Validation of the general routine}\label{validation-of-the-general-routine}

\textbf{The Woodbury routine is validated against a direct GLS computation and a seeded Monte Carlo benchmark that both exercise the common-close path.}

\begin{itemize}\setlength{\itemsep}{2pt}\setlength{\parskip}{0pt}
\item A direct computation assembles and inverts $\Sigma$ without the Woodbury identity over a grid of designs including $J_2 > 0$.
\item Agreement with the Woodbury routine is at the level of numerical round-off.
\item A seeded Monte Carlo at $(J_0, J_1, J_2) = (2, 2, 2)$ confirms that the empirical variance of $\hat\gamma$ matches the analytic value.
\end{itemize}

\subsection{Application: Alzheimer's disease imaging trial}\label{application-alzheimers-disease-imaging-trial}

\textbf{The Alzheimer's application uses MIRIAD-derived parameters with the correctly combined within-visit residual, which is essential to reproduce Frost's reported sample sizes.}

\begin{itemize}\setlength{\itemsep}{2pt}\setlength{\parskip}{0pt}
\item The slope variance and combined within-visit residual are taken from Frost's MIRIAD estimates.
\item The example assumes a one-year observation interval and a treatment effect of one quarter.
\item Using only the small measurement-error term understates the residual by two orders of magnitude and yields degenerate sample sizes.
\end{itemize}

\section{Discussion}\label{discussion}

\textbf{The work formalizes and extends Frost's framework with explicit variance formulas for run-in and common-close designs.}

\begin{itemize}\setlength{\itemsep}{2pt}\setlength{\parskip}{0pt}
\item The results provide explicit treatment-effect variance formulas.
\item They cover designs incorporating both run-in and common-close observations.
\item Three substantive points warrant further discussion.
\end{itemize}

\textbf{The benefit of run-in observations is governed by the slope-to-noise ratio and is largest for outcomes such as neuroimaging.}

\begin{itemize}\setlength{\itemsep}{2pt}\setlength{\parskip}{0pt}
\item The efficiency gain depends on the ratio of slope variance to residual variance.
\item When that ratio is large, even one or two run-in observations cut sample sizes substantially.
\item For cognitive outcomes with a less favorable ratio, gains are modest and must be weighed against longer study duration.
\end{itemize}

\textbf{Common-close observations sharpen the treatment-effect estimator at low cost, with diminishing returns beyond the first few visits.}

\begin{itemize}\setlength{\itemsep}{2pt}\setlength{\parskip}{0pt}
\item They inform individual trajectories without keeping participants on treatment.
\item Both groups contribute to slope-variability estimation once treatment has ceased.
\item Marginal benefit diminishes, but the first one or two visits give meaningful reductions when the slope-to-noise ratio is large.
\end{itemize}

\textbf{The derived formulas are consistent with Wang's two-period model while making the three phases explicit and separately optimizable.}

\begin{itemize}\setlength{\itemsep}{2pt}\setlength{\parskip}{0pt}
\item Wang's common-slope scenario corresponds to setting the run-in and placebo slopes equal.
\item Their different-slope scenario corresponds to the general case treated here.
\item The present formulation separates run-in, treatment, and common-close phases, enabling per-phase optimization.
\end{itemize}

\textbf{The derivations rest on simplifying assumptions whose relaxation is feasible but typically forfeits closed-form simplicity.}

\begin{itemize}\setlength{\itemsep}{2pt}\setlength{\parskip}{0pt}
\item They assume equally spaced observations, linear trajectories, and a random intercept-plus-slope structure.
\item Extensions to unequal spacing, nonlinear progression, or richer random effects remain within the Woodbury framework but lose simple closed forms.
\item No dropout is assumed, though the Dawson-Lagakos pattern-mixture approach can adjust for anticipated attrition.
\end{itemize}

\section{Data availability statement}\label{data-availability-statement}

\textbf{All materials supporting the paper are openly available in the project repository.}

\begin{itemize}\setlength{\itemsep}{2pt}\setlength{\parskip}{0pt}
\item These include the R package, the four companion manuscripts, and the Mathematica derivation scripts.
\item The shared bibliography is likewise available.
\item Specific paths are enumerated in the Reproducibility section.
\end{itemize}

\section{Reproducibility}\label{reproducibility}

\textbf{The manuscript build relies on the in-package source together with a small set of testing and rendering packages.}

\begin{itemize}\setlength{\itemsep}{2pt}\setlength{\parskip}{0pt}
\item Rendering used R version 4.4.x.
\item The principal dependencies are the in-package source plus tinytest, rmarkdown, and bookdown.
\item Supporting Mathematica symbolic derivations reside in the analysis scripts directory.
\end{itemize}

\textbf{The variance results are deterministic, and the Monte Carlo benchmark reported in the validation subsection uses a fixed seed.}

\begin{itemize}\setlength{\itemsep}{2pt}\setlength{\parskip}{0pt}
\item Given fixed input parameters the formulas exercise no random-number generation.
\item The validation subsection compares the Woodbury routine with a direct GLS computation over a grid that includes common-close configurations.
\item Its Monte Carlo benchmark at $(J_0, J_1, J_2) = (2, 2, 2)$ sets the seed 20260418 at its entry point.
\end{itemize}

\textbf{The repository organizes the package source, derivation scripts, companion manuscripts, and bibliography into well-defined locations.}

\begin{itemize}\setlength{\itemsep}{2pt}\setlength{\parskip}{0pt}
\item The R package source and the Mathematica scripts each have a dedicated directory.
\item The three companion manuscripts share a report subtree.
\item This manuscript source and the shared bibliography sit in the run-in power report directory.
\end{itemize}

\end{document}
